\documentclass[a4paper,11pt]{article}
\usepackage[utf8]{inputenc}
\usepackage[T1]{fontenc}
\usepackage[french]{babel}
\usepackage{lmodern}
\usepackage{amsmath,amssymb,amsthm}
\usepackage{mathrsfs}
\usepackage{enumitem}
\usepackage{multicol}

\setlist[enumerate]{label=\textbf{\textcolor{Couleur8}{\arabic*.}}, left=1em}

% Si vous voulez aussi modifier les listes enumerate imbriquées :
\setlist[enumerate,2]{label=\textcolor{Couleur8}{\alph*)}}
\setlist[enumerate,3]{label=\textcolor{Couleur8}{\roman*)}}
\setlist[enumerate,4]{label=\textcolor{Couleur8}{\Alph*)}}
\usepackage{graphicx}
\usepackage{xcolor}
\usepackage{tcolorbox}
\usepackage{geometry}
\usepackage{fancyhdr}
\usepackage{titlesec}
\usepackage{cancel}
\usepackage{ifthen}
\usepackage{tikz}
\usetikzlibrary{arrows.meta}
\usepackage{tkz-tab}
\usepackage{eurosym}

% OPTION PROF/ÉLÈVE
% Décommentez UNE des deux lignes suivantes :
\newboolean{prof}\setboolean{prof}{true}   % Version professeur (avec corrections des devoirs)
\newboolean{prof}\setboolean{prof}{true}  % Version élève (sans corrections des devoirs)

\definecolor{Couleur1}{HTML}{4A5D7A} %Th
\definecolor{Couleur2}{HTML}{A5B5D6} %Prop
\definecolor{Couleur3}{HTML}{A8C68A} %Méthode
\definecolor{Couleur6}{HTML}{8A9CC1} %Def
\definecolor{Couleur7}{HTML}{E6B459} %Rq Voc
\definecolor{Couleur8}{HTML}{D36740} %Titres
\definecolor{correction}{HTML}{D36740} % Couleur pour les corrections
\definecolor{coopmathscorr}{HTML}{F15929} % Couleur corrections coopmaths

% Configuration des marges
\geometry{
	top=2cm,
	bottom=2cm,
	inner=1.5cm,
	outer=1.5cm,
}

% Police
\renewcommand{\familydefault}{\sfdefault}

% Compteurs
\newcounter{seancenum}
\newcounter{seancenumD}
\setcounter{seancenum}{1}
\setcounter{seancenumD}{1}

% Configuration des en-têtes et pieds de page
\pagestyle{fancy}
\fancyhf{}
\ifthenelse{\boolean{prof}}{
    \fancyhead[L]{\textcolor{Couleur8}{\textbf{Corrections Automatismes - Classe de seconde - Version Professeur}}}
}{
    \fancyhead[L]{\textcolor{Couleur8}{\textbf{Corrections Devoirs - Séance 18 - Classe de seconde}}}
}
\fancyhead[R]{\textcolor{Couleur8}{\textbf{2025-2026}}}
\fancyfoot[L]{\textcolor{Couleur8}{Mme Bertail et Mme Pinto}}
\fancyfoot[R]{\textcolor{Couleur8}{\thepage}}
\renewcommand{\headrulewidth}{1pt}
\renewcommand{\headrule}{\hbox to\headwidth{\color{Couleur8}\leaders\hrule height \headrulewidth\hfill}}

% Environnements personnalisés
\newtcolorbox{seance}[1]{
    colback=Couleur6!10,
    colframe=Couleur1!65,
    fonttitle=\bfseries\large,
    title=Correction Séance \theseancenum #1,
    left=5mm,
    right=5mm,
    top=3mm,
    bottom=3mm,
    arc=0mm,
    after=\stepcounter{seancenum}
}

\newtcolorbox{seanceD}[1]{
    colback=Couleur6!5,
    colframe=Couleur6!45,
    fonttitle=\bfseries\large,
    title=Correction Devoirs - Séance \theseancenumD #1,
    left=5mm,
    right=5mm,
    top=3mm,
    bottom=3mm,
    arc=0mm,
    after=\stepcounter{seancenumD}
}

\begin{document}

\setcounter{seancenumD}{18}
\newpage
\begin{seanceD}{} %18

\begin{enumerate}
\item On écrit les tailles en écriture scientifique pour les comparer :\\
- Cellule 1 : $39{,}18\times 10^{-7} = 3{,}918\times 10^{-6}$ mm\\
- Cellule 2 : $470{,}16\times 10^{-8} = 4{,}7016\times 10^{-6}$ mm\\
- Cellule 3 : $322{,}02\times 10^{-6} = 3{,}2202\times 10^{-4}$ mm\\
- Cellule 4 : $505{,}2\times 10^{-9} = 5{,}052\times 10^{-7}$ mm\\

On a donc : $3{,}2202\times 10^{-4} > 4{,}7016\times 10^{-6} > 3{,}918\times 10^{-6} > 5{,}052\times 10^{-7}$\\
Donc il s'agit de la {\bfseries \color{correction}Cellule 3} qui a la taille la plus importante.\\
La bonne réponse est la réponse {\bfseries \color{correction}B}.

\item Le coefficient multiplicateur associé à une baisse de $5\,\%$ est $1-0{,}05=0{,}95$.\\
Le coefficient multiplicateur réciproque est donc $\dfrac{1}{0{,}95}$.\\
On en déduit que le taux réciproque est $\dfrac{1}{0{,}95}-1$ ou $\dfrac{0{,}05}{0{,}95}$.\\
Le taux à appliquer pour que cet article retrouve son prix initial est donné par ${\color{correction}\boldsymbol{\dfrac{0{,}05}{0{,}95}}}$.\\
La bonne réponse est la réponse {\bfseries \color{correction}D}.

\item $(20{,}1)^2>(20,05)^2$ car $20,10 > 20,05$

\item On peut simplifier l'expression :\\
$\begin{aligned}
\dfrac{8x^{3}}{\dfrac{2}{x^5}}&=8x^{3} \times \dfrac{x^5}{2}\\
&=\dfrac{8x^{8}}{2}\\
&=4x^{8}.
\end{aligned}$\\
La bonne réponse est la réponse {\bfseries \color{correction}A}.

\item On a :\\
$\begin{aligned}
A&=\dfrac{7}{10}+\dfrac{45}{100}\\
&=0{,}7+0{,}45\\
&={\color{correction}\boldsymbol{1{,}15}}
\end{aligned}$\\
La bonne réponse est la réponse {\bfseries \color{correction}A}.

\end{enumerate}

\end{seanceD}

\end{document}